%%
%% Springer Nature LaTeX Template for Quantum Information Processing
%% Converted from arXiv submission format
%%
\documentclass[pdflatex,sn-mathphys-num]{sn-jnl}

\usepackage{graphicx}
\usepackage{booktabs}
\usepackage{amsmath,amssymb,amsfonts}
\usepackage{amsthm}
\usepackage{algorithm}
\usepackage{algpseudocode}

\raggedbottom

\begin{document}

\title[Quantum PN Neuron Architecture for EEG Analysis]{Quantum Positive-Negative Neuron Architecture for Multi-Channel EEG Analysis: Hardware Validation and Empirical Limits}

\author*[1]{\fnm{James} \sur{Appenzeller}}

\affil[1]{\orgname{Independent Researcher}, \city{West Windsor}, \state{NJ}, \country{USA}}

\abstract{Hardware validation on IBM's 133-qubit Heron processor demonstrates that a quantum neural architecture can reliably distinguish neural configurations with 99.3\% discrimination quality. The A-Gate circuit encodes excitatory-inhibitory dynamics from EEG signals using paired qubits, achieving $O(M)$ gate complexity for $M$-channel correlation encoding versus $O(M^2)$ classical computation on the symmetric positive definite (SPD) manifold. Transpilation to native CZ gates confirms linear scaling ($14.1M - 17.5$, $R^2 > 0.99$). Evaluation on the CHB-MIT Scalp EEG Database using leave-one-subject-out (LOSO) cross-validation across 22 subjects reveals a critical encoding bottleneck: the quantum fidelity classifier achieves 53.4\% AUC compared to 62.5\% for an 8-channel classical baseline. Notably, per-subject AUC ranges from 0.183 to 1.000, with sub-chance performance in some cases suggesting what we term \emph{manifold polarity}---patient-specific geometric orientations that may cause cross-subject template mismatch. The gap between quantum and classical performance indicates that mapping EEG features to circuit parameters, not the A-Gate architecture itself, limits classification. We provide rigorous complexity analysis with corrections to common quantum advantage overclaims, establishing both the promise and present limitations of quantum neural signal processing.}

\keywords{quantum computing, neural networks, EEG analysis, positive-negative neuron, hardware validation, seizure prediction}

\maketitle

\section{Introduction}
\label{sec:intro}

Seizure prediction from electroencephalogram (EEG) recordings remains a challenging problem in computational neuroscience \cite{rasheed2021}. Pre-ictal states exhibit complex spatiotemporal dynamics characterized by increased synchronization across cortical regions, cross-frequency coupling, and subtle phase relationships that emerge minutes to hours before seizure onset \cite{mormann2007}. The clinical significance of reliable seizure prediction is substantial: accurate prediction would enable timely intervention for the approximately 50 million people worldwide affected by epilepsy \cite{who2019}. Recent advances in deep learning have achieved promising results \cite{daoud2019}, yet fundamental limitations in classical correlation analysis motivate exploration of alternative computational paradigms.

Classical approaches to multi-channel EEG analysis require explicit computation of pairwise correlations. For a standard clinical 10-20 montage with 19 electrodes, this necessitates calculating 171 unique channel pairs. As EEG systems evolve toward higher spatial resolution with 64, 128, or 256 channels, this quadratic scaling presents an increasingly significant computational burden.

The Positive-Negative (PN) neuron model \cite{gupta2003} provides a biologically-inspired framework capturing excitatory-inhibitory (E-I) dynamics fundamental to neural computation. In biological circuits, E-I balance determines network stability, and disruption of this balance is a hallmark of epileptogenic tissue \cite{buzsaki2012}.

Quantum machine learning offers a potentially transformative approach to such high-dimensional correlation problems \cite{biamonte2017}. Current noisy intermediate-scale quantum (NISQ) devices \cite{preskill2018} enable exploration of hybrid quantum-classical algorithms \cite{cerezo2021}, including quantum neural network architectures \cite{farhi2018,abbas2021}. Recent work has explored quantum approaches to EEG analysis \cite{gandhi2014,chen2024qeegnet,olvera2024}, including hybrid quantum-classical architectures and quantum-enhanced feature extraction. However, none have validated on quantum hardware or systematically investigated the encoding bottleneck.

This paper proposes a quantum implementation of the PN neuron architecture that leverages superposition and entanglement to encode E-I dynamics and inter-channel correlations more efficiently than classical methods. Unlike prior theoretical work, we provide empirical validation on both clinical EEG data and IBM quantum hardware.

The primary contributions are:
\begin{enumerate}
    \item A quantum circuit architecture (the ``A-Gate'') mapping PN neuron parameters to qubit rotations with $O(M)$ gate complexity.
    \item Hardware validation on IBM's 133-qubit Heron processor demonstrating 99.3\% configuration discrimination and confirming theoretical scaling.
    \item Rigorous evaluation on the CHB-MIT dataset with leave-one-subject-out cross-validation, establishing performance bounds and identifying the encoding bottleneck.
    \item Explicit corrections to common overclaims regarding quantum advantage in machine learning.
\end{enumerate}

\section{Background}
\label{sec:background}

\subsection{The Positive-Negative Neuron Model}
\label{sec:pn-model}

The PN neuron model describes neural dynamics through coupled differential equations governing excitatory and inhibitory state variables. For each neural unit, three parameters are maintained: amplitude $a$ governing excitatory activation magnitude, phase $b$ encoding temporal dynamics and inter-unit coupling, and coupling strength $c$ characterizing inhibitory influence.

The temporal evolution follows first-order dynamics \cite{gupta2003}:
\begin{equation}
\frac{da}{dt} = -\lambda_a \cdot a + f(t)(1 - a)
\end{equation}
\begin{equation}
\frac{dc}{dt} = +\lambda_c \cdot c + f(t)(1 - c)
\end{equation}
where $f(t)$ represents the normalized input signal and $\lambda_a$, $\lambda_c$ are decay constants. The asymmetric structure---negative decay for excitation and positive accumulation for inhibition---reflects biological differences between these neural processes.

\subsection{Quantum Computing Preliminaries}
\label{sec:quantum-prelim}

A quantum system of $n$ qubits exists in a $2^n$-dimensional complex Hilbert space. The state $|\psi\rangle$ is described by $2^n$ complex amplitudes, with measurement probability $|\alpha_i|^2$ for basis state $|i\rangle$. While this exponential state space provides the foundation for quantum advantage, measurement collapses the superposition to a single classical outcome---a distinction crucial for honest assessment of quantum advantage claims.

Relevant quantum operations include single-qubit gates (Hadamard H, phase P$(\theta)$, rotations $R_x$, $R_y$, $R_z$) and two-qubit gates (controlled rotations CR$_y$, CR$_z$, CNOT, CZ). The CZ gate is particularly relevant as it constitutes the native two-qubit gate on IBM's Heron processors.

With these foundations established, we now describe the quantum circuit architecture that maps PN neuron dynamics onto qubit operations.

\section{Quantum PN Neuron Architecture}
\label{sec:architecture}

\subsection{Single-Channel Encoding: The A-Gate}
\label{sec:agate}

For each EEG channel, the quantum PN architecture allocates two qubits representing excitatory (E) and inhibitory (I) components. The PN parameters $(a, b, c)$ are encoded through a two-layer circuit termed the ``A-Gate.''

The first layer applies per-qubit encoding using an H-P-R-P-H structure. For the excitatory qubit, this sequence consists of H, P$(b)$, $R_x(2a)$, P$(b)$, H. The inhibitory qubit follows the same structure but substitutes $R_y(2c)$ for the central rotation. The shared phase parameter $b$ appears on all four phase gates, encoding temporal coupling through quantum phase.

The second layer establishes bidirectional E-I coupling through controlled rotation gates: CR$_y(\pi/4)$ with excitatory control and CR$_z(\pi/4)$ with inhibitory control. This bidirectional structure reflects biological mutual regulation between excitation and inhibition.

The complete single-channel encoding requires 12 gates with circuit depth 7: four Hadamard gates (two per qubit) establish and collapse superposition, four phase gates P$(b)$ encode temporal coupling through the shared phase parameter, two parameterized rotations ($R_x(2a)$ on the excitatory qubit and $R_y(2c)$ on the inhibitory qubit) encode activation magnitudes, and two controlled rotations (CR$_y$ and CR$_z$) establish bidirectional E-I coupling.

\begin{figure}[htbp]
\centering
\includegraphics[width=0.85\textwidth]{figures/agate_circuit.png}
\caption{The A-Gate circuit encodes PN neuron parameters $(a, b, c)$ using two qubits. The excitatory qubit undergoes H-P$(b)$-$R_x(2a)$-P$(b)$-H encoding while the inhibitory qubit uses $R_y(2c)$. Bidirectional E-I coupling is established through CR$_y(\pi/4)$ and CR$_z(\pi/4)$ gates.}
\label{fig:agate}
\end{figure}

\subsection{Multi-Channel Entanglement Architecture}
\label{sec:multichannel}

For $M$ channels, the architecture employs $2M$ qubits plus one ancilla, totaling $2M+1$ qubits. Inter-channel correlations are captured through two complementary strategies:

\textbf{Ring topology}: Sequential CNOT gates connect adjacent channel qubits within each manifold: $E_1 \rightarrow E_2 \rightarrow \cdots \rightarrow E_M$ and $I_1 \rightarrow I_2 \rightarrow \cdots \rightarrow I_M$. This requires $2(M-1)$ CNOT gates and encodes nearest-neighbor correlations corresponding to spatial electrode adjacency.

\textbf{Global ancilla coupling}: An ancilla qubit initialized in Hadamard superposition connects to all excitatory qubits via CZ gates. This requires $M$ additional CZ gates and enables detection of widespread synchronization patterns characteristic of pre-ictal states.

\subsection{Gate Count Analysis}
\label{sec:gate-count}

The total gate requirement for $M$ channels can be enumerated as follows. Each channel requires 12 gates for A-Gate encoding, contributing $12M$ gates. The ring topology connecting adjacent qubits within each manifold requires $M-1$ CNOT gates for the excitatory chain and $M-1$ for the inhibitory chain, totaling $2(M-1) = 2M-2$ gates. The ancilla qubit requires one Hadamard gate for initialization, and global coupling adds $M$ CZ gates connecting the ancilla to all excitatory qubits. Summing these contributions:
\begin{equation}
\text{Total gates} = 12M + (2M - 2) + M + 1 = 15M - 1 = O(M)
\end{equation}

This linear scaling represents a fundamental improvement over $O(M^2)$ classical correlation computation. The linear gate count established above invites formal comparison with classical computational requirements, which we undertake in the following section.

\begin{figure}[htbp]
\centering
\includegraphics[width=0.9\textwidth]{figures/multichannel_circuit.png}
\caption{Multi-channel circuit architecture for $M=4$ channels (9 qubits). Each channel is encoded by an A-Gate, followed by ring topology entanglement within E and I manifolds, and global ancilla coupling via CZ gates. Total gate count: $15M - 1 = O(M)$.}
\label{fig:multichannel}
\end{figure}

\section{Complexity Analysis}
\label{sec:complexity}

\subsection{Classical Baseline}
\label{sec:classical-complexity}

Classical analysis of $M$-channel EEG data with $T$ time samples involves several computational phases. Feature extraction for each channel independently requires $O(M \cdot T)$ operations. Pairwise correlation computation---specifically, phase locking value (PLV) calculation across $M(M-1)/2$ unique channel pairs with $T$ time samples each---contributes $O(M^2 \cdot T)$ operations and dominates the computational cost. Finally, template matching against $M^2$-dimensional feature vectors that encode these pairwise correlations adds $O(M^2)$ operations per comparison. The aggregate classical complexity is therefore $O(M^2 \cdot T)$ for preprocessing and $O(M^2)$ per classification.

\subsection{Quantum Complexity}
\label{sec:quantum-complexity}

The quantum approach restructures these computations to achieve linear scaling in $M$. Classical preprocessing to extract PN parameters from each channel remains $O(M \cdot T)$---this cost is unchanged from the classical baseline. However, quantum encoding requires only $O(M)$ gates, capturing all pairwise correlations implicitly through entanglement without explicit enumeration of channel pairs. Template matching via the SWAP test \cite{buhrman2001} requires $O(M)$ additional gates rather than $O(M^2)$ classical operations. The total quantum complexity is therefore $O(M \cdot T)$ for preprocessing plus $O(M)$ for encoding and classification, eliminating the quadratic bottleneck.

\begin{table}[htbp]
\centering
\caption{Complexity Comparison: Classical versus Quantum Approaches}
\label{tab:complexity}
\begin{tabular}{lccc}
\toprule
Operation & Classical & Quantum & Advantage \\
\midrule
Correlation encoding & $O(M^2)$ & $O(M)$ & $M\times$ \\
Template matching & $O(M^2)$ & $O(M)$ & $M\times$ \\
Parameter storage & $O(M^2)$ & $O(M)$ & $M\times$ \\
\bottomrule
\end{tabular}
\end{table}

\subsection{Clarifications on Quantum Information Capacity}
\label{sec:holevo}

A persistent overclaim in quantum machine learning conflates Hilbert space dimensionality with information capacity. While $2M$ qubits span a $2^{2M}$-dimensional space, the Holevo bound \cite{holevo1973} limits extractable classical information to at most $2M$ bits per measurement.

The quantum advantage arises not from raw information throughput but from how interference and entanglement process correlations during computation. The structure of entangled states encodes correlation information that requires only $O(M)$ resources to prepare, even though classical representation would require $O(M^2)$ resources.

\begin{figure}[htbp]
\centering
\includegraphics[width=0.72\textwidth]{figures/complexity_comparison.png}
\caption{Computational complexity scaling: classical $O(M^2)$ versus quantum $O(M)$ for correlation encoding. The advantage factor grows linearly with channel count, reaching 8$\times$ for our 8-channel implementation and approaching 64$\times$ for high-density 64-channel systems.}
\label{fig:complexity}
\end{figure}

\subsection{Measurement Statistics Overhead}
\label{sec:measurement}

The SWAP test yields fidelity $F = |\langle\psi|\phi\rangle|^2$ through repeated measurements. Estimating this probability to precision $\varepsilon$ requires $O(1/\varepsilon^2)$ shots by the Chernoff bound. For $\varepsilon = 0.01$, approximately 10,000 shots are required. This overhead does not negate the quantum advantage for correlation encoding but introduces a multiplicative constant for practical implementations. These theoretical scaling properties motivate empirical evaluation on clinical EEG data and quantum hardware.

\section{Experimental Setup}
\label{sec:setup}

\subsection{CHB-MIT Scalp EEG Database}
\label{sec:dataset}

Evaluation was performed on the CHB-MIT Scalp EEG Database \cite{goldberger2000,shoeb2009}, a standard benchmark for seizure prediction research. The database contains continuous EEG recordings from 22 pediatric patients with intractable seizures, collected at Boston Children's Hospital.

\begin{table}[htbp]
\centering
\caption{CHB-MIT Dataset Characteristics}
\label{tab:chbmit}
\begin{tabular}{ll}
\toprule
Parameter & Value \\
\midrule
Subjects & 22 \\
Recording montage & 10-20 international system \\
Available channels & 23 bipolar derivations \\
Sampling rate & 256 Hz \\
Total segments & 504 \\
Ictal/preictal segments & 284 (142 each) \\
Interictal segments & 220 \\
Segment duration & 30 seconds \\
\bottomrule
\end{tabular}
\end{table}

\subsection{Channel Selection}
\label{sec:channels}

Due to quantum simulation constraints (classical simulation limits of $\sim$30 qubits), we selected eight channels covering bilateral frontal-temporal regions: FP1-F7, F7-T7, FP1-F3, and F3-C3 from the left hemisphere, and FP2-F8, F8-T8, FP2-F4, and F4-C4 from the right hemisphere. This selection captures the primary regions involved in seizure propagation while maintaining computational tractability (17 qubits for the quantum circuit).

\subsection{Evaluation Protocol}
\label{sec:protocol}

We employed leave-one-subject-out (LOSO) cross-validation, the gold standard for seizure prediction evaluation that tests generalization across patients. For each of the 22 subjects, the classifier was trained on segments from the remaining 21 subjects and tested on the held-out subject's data. Predictions were then aggregated across all folds to compute final performance metrics. This protocol avoids within-subject data leakage and provides realistic estimates of clinical deployment performance.

\subsection{Encoding Strategies}
\label{sec:encoding}

Three encoding strategies were evaluated to map EEG signals to PN parameters, each representing a different hypothesis about which signal features best capture pre-ictal dynamics.

\textbf{V1 (PN Dynamics)}: Direct numerical integration of the PN differential equations with $\lambda_a = 0.1$, $\lambda_c = 0.05$, $dt = 0.01$. This approach preserves the original biological interpretation of the PN model, where the signal amplitude directly drives excitatory-inhibitory state evolution.

\textbf{V2 (Band Power)}: Spectral decomposition into standard EEG bands ($\delta$, $\theta$, $\alpha$, $\beta$, $\gamma$) with relative power mapping to PN parameters. This encoding leverages the neurophysiological correspondence between frequency bands and cortical processes: beta and gamma activity correlate with excitatory cortical processes, while delta and theta rhythms are associated with inhibitory and thalamocortical dynamics.

\textbf{V3 (PLV/Hilbert)}: Hilbert transform for instantaneous phase extraction with phase locking value computation between band-filtered signals. This encoding targets phase synchronization directly, motivated by evidence that pre-ictal states exhibit characteristic changes in inter-regional phase coupling.

\subsection{Classical Baseline}
\label{sec:baseline}

For fair comparison, we implemented an XGBoost \cite{chen2016xgboost} classifier with identical channel selection (8 channels) and additional evaluation with the full 18-channel montage:

\begin{table}[htbp]
\centering
\caption{Classical Baseline Configuration}
\label{tab:classical}
\begin{tabular}{lcc}
\toprule
Parameter & 8-Channel & 18-Channel \\
\midrule
Classifier & XGBoost & XGBoost \\
Estimators & 200 & 200 \\
Max depth & 6 & 6 \\
Bagging runs & 5 & 5 \\
Features & 816 & 2016 \\
\bottomrule
\end{tabular}
\end{table}

Features include spectral band powers, statistical moments, Hjorth parameters, and pairwise correlation coefficients. With this experimental framework in place, we first validate the circuit architecture on quantum hardware before presenting classification results.

\section{Hardware Validation}
\label{sec:hardware}

\subsection{IBM Quantum Platform}
\label{sec:ibm}

Hardware validation was performed on IBM's ibm\_torino backend, a 133-qubit Heron processor representing current state-of-the-art superconducting quantum hardware. Circuits were transpiled to the native gate set using Qiskit \cite{javadi2024} optimization level 3.

\begin{table}[htbp]
\centering
\begin{minipage}{0.42\textwidth}
\centering
\caption{IBM Quantum Hardware Specifications}
\label{tab:hardware}
\begin{tabular}{ll}
\toprule
Parameter & Value \\
\midrule
Backend & ibm\_torino \\
Processor & IBM Heron r1 \\
Total qubits & 133 \\
Native gates & CZ, RZ, SX, X \\
Median CZ error & $\sim$0.5\% \\
Median T1 & $\sim$150 $\mu$s \\
Median T2 & $\sim$200 $\mu$s \\
\bottomrule
\end{tabular}
\end{minipage}
\hfill
\begin{minipage}{0.54\textwidth}
\centering
\caption{Transpilation Results (8 Channels)}
\label{tab:transpile}
\begin{tabular}{lcc}
\toprule
Gate Type & Original & Transpiled \\
\midrule
H & 34 & --- \\
P & 32 & --- \\
CX/CNOT & 14 & --- \\
$R_x$, $R_y$ & 16 & --- \\
CR$_y$, CR$_z$ & 16 & --- \\
CZ & 8 & 97 \\
SX & --- & 204 \\
RZ & --- & 152 \\
X & --- & 10 \\
\midrule
Total & 120 & 481 \\
Depth & 17 & 114 \\
\bottomrule
\end{tabular}
\end{minipage}
\end{table}

\subsection{Transpilation Analysis}
\label{sec:transpile}

The transpiled circuit uses 97 native CZ gates, closely matching the theoretical prediction of $14.1 \times 8 - 17.5 \approx 95$ two-qubit gates.

\subsection{Scaling Experiment}
\label{sec:scaling}

To verify $O(M)$ scaling, circuits were generated for $M \in \{2, 4, 6, 8\}$ channels and transpiled:

\begin{table}[htbp]
\centering
\caption{Two-Qubit Gate Scaling}
\label{tab:scaling}
\begin{tabular}{ccccc}
\toprule
Channels ($M$) & Qubits & Original Depth & Transpiled Depth & CZ Gates \\
\midrule
2 & 5 & 12 & 33 & 14 \\
4 & 9 & 14 & 62 & 34 \\
6 & 13 & 16 & 73 & 67 \\
8 & 17 & 18 & 114 & 97 \\
\bottomrule
\end{tabular}
\end{table}

Linear regression yields:
\begin{equation}
\text{CZ gates} = 14.1 \cdot M - 17.5 \quad (R^2 > 0.99)
\end{equation}

This confirms $O(M)$ scaling in hardware-native gate counts, validating the theoretical complexity analysis.

\begin{figure}[htbp]
\centering
\includegraphics[width=0.75\textwidth]{figures/ibm_scaling.png}
\caption{Hardware scaling validation on ibm\_torino. Two-qubit (CZ) gate counts follow the linear relationship $14.1M - 17.5$ with $R^2 > 0.99$, confirming $O(M)$ complexity.}
\label{fig:scaling}
\end{figure}

\subsection{Configuration Discrimination}
\label{sec:discrimination}

Five distinct neural configurations were tested to assess the architecture's ability to distinguish different EEG states: synchronized (all channels with aligned phase, $b = 0$), desynchronized (random phase distribution across channels), half-synchronized (mixed alignment), excitatory-dominant (elevated $a$ parameters representing high excitation), and inhibitory-dominant (elevated $c$ parameters representing high inhibition). Each configuration was executed with 8192 shots. Discrimination quality was computed as $1 - \bar{F}_{\text{off-diag}}$ where $\bar{F}_{\text{off-diag}}$ is the mean off-diagonal Hellinger fidelity \cite{fuchs1999}.

\begin{table}[htbp]
\centering
\caption{Configuration Discrimination Quality}
\label{tab:discrimination}
\begin{tabular}{lcc}
\toprule
Execution Mode & Off-Diagonal Mean & Discrimination Quality \\
\midrule
Ideal (statevector) & 0.0182 & 98.2\% \\
Noisy simulation & 0.0089 & 99.1\% \\
Hardware (ibm\_torino) & 0.0065 & 99.3\% \\
\bottomrule
\end{tabular}
\end{table}

The hardware achieves 99.3\% discrimination quality, demonstrating that the quantum circuit reliably produces distinguishable output distributions for different neural configurations.

\begin{figure}[htbp]
\centering
\includegraphics[width=0.7\textwidth]{figures/ibm_discrimination_heatmap.png}
\caption{Configuration discrimination matrix from ibm\_torino hardware execution. Diagonal elements are 1.0 (self-similarity); off-diagonal values show low cross-configuration Hellinger fidelity, yielding 99.3\% discrimination quality.}
\label{fig:discrimination}
\end{figure}

\subsection{Fidelity Analysis}
\label{sec:fidelity}

Hellinger fidelity between execution modes quantifies hardware accuracy:

\begin{table}[htbp]
\centering
\caption{Hellinger Fidelity by Configuration}
\label{tab:fidelity}
\begin{tabular}{lccc}
\toprule
Configuration & Ideal vs Noisy & Ideal vs Hardware & Noisy vs Hardware \\
\midrule
synchronized & 0.038 & 0.021 & 0.015 \\
desynchronized & 0.323 & 0.242 & 0.209 \\
half\_sync & 0.175 & 0.122 & 0.098 \\
excitatory & 0.077 & 0.045 & 0.034 \\
inhibitory & 0.101 & 0.052 & 0.036 \\
\midrule
Mean & 0.143 & 0.096 & 0.078 \\
\bottomrule
\end{tabular}
\end{table}

Hardware execution shows closer agreement with noisy simulation than with ideal statevector computation, as expected. The ``desynchronized'' configuration shows highest deviation, likely due to greater sensitivity to gate errors when phase relationships are randomized. Having confirmed that the circuit operates correctly on quantum hardware, we now evaluate its classification performance on clinical EEG data.

\section{Results}
\label{sec:results}

\subsection{Encoding Ablation Study}
\label{sec:ablation}

Table~\ref{tab:encoding} summarizes performance across encoding strategies:

\begin{table}[htbp]
\centering
\caption{Encoding Strategy Comparison (22-Subject LOSO)}
\label{tab:encoding}
\begin{tabular}{lccccc}
\toprule
Encoding & AUC & Accuracy & Precision & Recall & F1 \\
\midrule
V1 (PN dynamics) & 0.444 & 0.593 & 0.581 & 1.000 & 0.735 \\
V2 (band power) & \textbf{0.534} & 0.550 & 0.606 & 0.574 & 0.590 \\
V3 (PLV/Hilbert) & 0.529 & 0.563 & 0.563 & 1.000 & 0.721 \\
\bottomrule
\end{tabular}
\end{table}

The V2 (band power) encoding achieves the highest AUC at 53.4\%, outperforming both direct PN dynamics integration (44.4\%) and phase-based encoding (52.9\%). However, all encoding strategies perform below chance-level discrimination when considering the class imbalance.

\begin{figure}[htbp]
\centering
\includegraphics[width=0.85\textwidth]{figures/encoding_ablation.png}
\caption{Encoding strategy comparison on CHB-MIT dataset (22-subject LOSO). V2 (band power) encoding achieves the best quantum performance at 53.4\% AUC, but remains below the 8-channel classical baseline (62.5\%) and 18-channel ceiling (81.98\%).}
\label{fig:encoding}
\end{figure}

\subsection{Classical Baseline Comparison}
\label{sec:comparison}

\begin{table}[htbp]
\centering
\caption{Quantum vs Classical Performance}
\label{tab:qvc}
\begin{tabular}{lcccc}
\toprule
Method & Channels & Features & AUC & Accuracy \\
\midrule
Quantum (V2) & 8 & 24 (PN params) & 0.534 & 0.550 \\
Classical (XGBoost) & 8 & 816 & 0.625 & 0.595 \\
Classical (XGBoost) & 18 & 2016 & \textbf{0.820} & \textbf{0.757} \\
\bottomrule
\end{tabular}
\end{table}

Three key findings emerge from this comparison. First, the 8-channel classical baseline outperforms the quantum approach (62.5\% vs 53.4\% AUC) despite using identical channel selection, indicating that the performance gap is not due to channel count. Second, the 18-channel classical baseline achieves 81.98\% AUC, establishing a performance ceiling with current feature engineering and demonstrating that additional channels substantially improve discrimination. Third, the quantum approach uses significantly fewer parameters (24 PN parameters versus 816 or 2016 classical features), suggesting that parameter efficiency is not the bottleneck.

\subsection{Per-Subject Analysis}
\label{sec:persubject}

Performance varies substantially across subjects, as is typical in seizure prediction:

\begin{table}[htbp]
\centering
\caption{Per-Subject AUC Range (Classical 8-Channel)}
\label{tab:persubject}
\begin{tabular}{lccc}
\toprule
Metric & Min & Mean & Max \\
\midrule
AUC & 0.183 (chb17) & 0.699 & 1.000 (chb22) \\
\bottomrule
\end{tabular}
\end{table}

Several subjects achieve excellent discrimination (chb01: 0.929, chb05: 0.920, chb18: 0.950, chb20: 0.988, chb22: 1.000), while others prove challenging (chb14: 0.344, chb17: 0.183). This heterogeneity reflects known variability in seizure prediction and motivates subject-specific adaptation strategies. These results point to a clear conclusion about where the bottleneck lies, which we examine in the following discussion.

\begin{figure}[htbp]
\centering
\includegraphics[width=0.95\textwidth]{figures/per_subject_auc.png}
\caption{Per-subject AUC distribution for the classical 8-channel baseline. Performance varies substantially across subjects, with several achieving near-perfect discrimination while others remain challenging. This heterogeneity is characteristic of seizure prediction tasks.}
\label{fig:subjects}
\end{figure}

\section{Discussion}
\label{sec:discussion}

\subsection{The Encoding Bottleneck}
\label{sec:bottleneck}

The results reveal a critical insight: the encoding strategy---mapping EEG signals to quantum circuit parameters---constitutes the primary performance bottleneck, not the quantum circuit architecture itself. This finding aligns with theoretical work showing that data encoding fundamentally determines the expressive power of variational quantum models \cite{schuld2021encoding}, and that encoded quantum states can concentrate toward uninformative distributions \cite{li2022concentration}.

Three lines of evidence support this conclusion. First, hardware validation demonstrates excellent discrimination (99.3\%) between distinct configurations, indicating that the circuit architecture can faithfully represent and distinguish different neural states---the quantum representation itself is not the limiting factor. Second, the V1 $\rightarrow$ V2 improvement (44.4\% $\rightarrow$ 53.4\% AUC) demonstrates that encoding refinement yields measurable gains, confirming that the encoding choice materially affects classification performance. Third, the persistent gap between quantum (53.4\%) and classical (62.5\%) performance using identical channels indicates systematic information loss during the quantum encoding process, rather than any limitation of the circuit's discriminative capacity.

The classical baseline benefits from 816 engineered features capturing spectral, statistical, and correlation information. Notably, the covariance matrices underlying these correlation features reside on the symmetric positive definite (SPD) manifold, a Riemannian structure whose geodesic distances have proven effective for EEG classification in the brain-computer interface literature. The quantum approach compresses this to 24 PN parameters (3 per channel), discarding potentially discriminative information---including the full SPD geometry.

\subsection{Implications for Quantum Advantage}
\label{sec:implications}

The complexity analysis demonstrates clear theoretical advantage: $O(M)$ quantum versus $O(M^2)$ classical for correlation encoding. Hardware validation confirms this scaling holds in practice (CZ gates $\sim 14.1M - 17.5$).

However, theoretical scaling advantage does not guarantee practical classification advantage. Unlike variational approaches susceptible to barren plateaus \cite{mcclean2018}, the A-Gate architecture uses physics-informed fixed parameters, avoiding trainability issues but requiring careful encoding design. The current encoding strategies fail to preserve sufficient discriminative information for seizure prediction. Achieving practical quantum advantage will require richer encoding strategies that preserve more signal information while remaining compatible with circuit depth constraints, larger channel counts ($M > 50$) where the $O(M)$ versus $O(M^2)$ difference becomes pronounced, and hybrid approaches that combine quantum correlation encoding with classical feature extraction to leverage the strengths of each paradigm.

\subsection{Hardware Readiness}
\label{sec:readiness}

The ibm\_torino results demonstrate that current quantum hardware can execute the 17-qubit, 97-CZ-gate circuits required for 8-channel analysis, produce distinguishable output distributions for different neural configurations, and maintain sufficient coherence through circuit depths of 114 transpiled layers. These capabilities were uncertain prior to hardware validation and represent meaningful progress toward practical quantum neural signal processing.

\subsection{Limitations}
\label{sec:limitations}

\textbf{Simulation constraints}: Classical simulation limits restricted evaluation to 8 channels. The theoretical advantage becomes more pronounced at higher channel counts inaccessible to current simulators.

\textbf{Encoding information loss}: The compression from raw EEG to PN parameters discards information that classical approaches retain through richer feature sets.

\textbf{Single dataset}: Evaluation on CHB-MIT alone limits generalizability claims. Validation on additional datasets (e.g., TUSZ, Bonn) would strengthen conclusions.

\textbf{Template-based classification}: The quantum fidelity approach uses fixed templates rather than learned decision boundaries, potentially limiting adaptability.

\textbf{Subject-specific geometry}: The substantial per-subject AUC variance---ranging from 0.183 to 1.000, with several subjects performing below chance---suggests that individual patients may occupy geometrically distinct regions of the feature space. We term this phenomenon \emph{manifold polarity}: the possibility that subjects' pre-ictal signatures project to opposing orientations on the SPD manifold, causing templates optimized for one group to anti-correlate with another. This interpretation remains speculative but warrants investigation, as it would have direct implications for encoding strategy design and subject-specific calibration.

\section{Conclusion}
\label{sec:conclusion}

This work presents the first comprehensive validation of a quantum neural architecture for EEG analysis, combining theoretical complexity analysis, clinical dataset evaluation, and quantum hardware execution.

Key contributions:
\begin{enumerate}
    \item The A-Gate architecture achieves $O(M)$ gate complexity for $M$-channel correlation encoding, confirmed by hardware measurements showing CZ gates $= 14.1M - 17.5$.
    \item IBM Heron processor validation demonstrates 99.3\% configuration discrimination, establishing hardware feasibility.
    \item CHB-MIT evaluation (22 subjects, LOSO) reveals the encoding bottleneck: quantum AUC of 53.4\% versus classical ceiling of 81.98\%.
    \item Honest assessment identifies encoding strategy---not circuit architecture---as the primary improvement target.
\end{enumerate}

The path forward requires developing encoding strategies that preserve discriminative EEG information while respecting quantum circuit constraints. The architecture is sound; the translation from neuroscience to quantum representation requires refinement. As quantum hardware continues to mature beyond the current NISQ era \cite{preskill2018}, architectures designed today with careful attention to both theoretical advantage and practical limitations will be positioned to deliver clinical impact.

\section*{Data Availability}
The CHB-MIT dataset is publicly available from PhysioNet \cite{goldberger2000}. IBM Quantum hardware results are archived in the project repository.

\section*{Code Availability}
Implementation code is available at \url{https://github.com/jappenzeller/QDNU}.

\section*{Author Contributions}
J.A. conceived the architecture, implemented the quantum circuits, conducted all experiments, and wrote the manuscript.

\section*{Competing Interests}
The author declares no competing interests.

\bibliography{references}

\begin{appendices}

\section{Gate Definitions}
\label{app:gates}

\subsection{Single-Qubit Gates}

\textbf{Hadamard Gate}:
\begin{equation}
H = \frac{1}{\sqrt{2}} \begin{pmatrix} 1 & 1 \\ 1 & -1 \end{pmatrix}
\end{equation}

\textbf{Phase Gate}:
\begin{equation}
P(\theta) = \begin{pmatrix} 1 & 0 \\ 0 & e^{i\theta} \end{pmatrix}
\end{equation}

\textbf{Rotation Gates}:
\begin{equation}
R_x(\theta) = \begin{pmatrix} \cos(\theta/2) & -i\sin(\theta/2) \\ -i\sin(\theta/2) & \cos(\theta/2) \end{pmatrix}
\end{equation}
\begin{equation}
R_y(\theta) = \begin{pmatrix} \cos(\theta/2) & -\sin(\theta/2) \\ \sin(\theta/2) & \cos(\theta/2) \end{pmatrix}
\end{equation}

\subsection{Two-Qubit Gates}

The controlled-Z (CZ) gate, native to IBM Heron processors:
\begin{equation}
CZ = \begin{pmatrix} 1 & 0 & 0 & 0 \\ 0 & 1 & 0 & 0 \\ 0 & 0 & 1 & 0 \\ 0 & 0 & 0 & -1 \end{pmatrix}
\end{equation}

\section{A-Gate Implementation}
\label{app:implementation}

The A-Gate circuit is implemented using Qiskit. The complete implementation is available in the project repository at \url{https://github.com/jappenzeller/QDNU}. The core circuit construction follows Algorithm~\ref{alg:agate}.

\begin{algorithm}[htbp]
\caption{A-Gate Circuit Construction}
\label{alg:agate}
\begin{algorithmic}[1]
\Require PN parameters $(a, b, c)$ where $a, c \in [0, 1]$ and $b \in [0, 2\pi]$
\Ensure Two-qubit quantum circuit encoding PN neuron state
\State Initialize 2-qubit circuit $\mathcal{C}$
\State \Comment{Excitatory qubit encoding}
\State Apply H to qubit 0
\State Apply P$(b)$ to qubit 0
\State Apply $R_x(2a)$ to qubit 0
\State Apply P$(b)$ to qubit 0
\State Apply H to qubit 0
\State \Comment{Inhibitory qubit encoding}
\State Apply H to qubit 1
\State Apply P$(b)$ to qubit 1
\State Apply $R_y(2c)$ to qubit 1
\State Apply P$(b)$ to qubit 1
\State Apply H to qubit 1
\State \Comment{E-I coupling}
\State Apply CR$_y(\pi/4)$ controlled by qubit 0, target qubit 1
\State Apply CR$_z(\pi/4)$ controlled by qubit 1, target qubit 0
\State \Return $\mathcal{C}$
\end{algorithmic}
\end{algorithm}

\section{Per-Subject Results}
\label{app:subjects}

\begin{table}[htbp]
\centering
\caption{Per-Subject AUC for Classical XGBoost Baseline (8 Channels). These results establish the classical performance ceiling for comparison with quantum approaches.}
\label{tab:subjects}
\small
\begin{tabular}{lclc}
\toprule
Subject & AUC & Subject & AUC \\
\midrule
chb01 & 0.929 & chb13 & 0.479 \\
chb02 & 0.783 & chb14 & 0.344 \\
chb03 & 0.550 & chb15 & 0.583 \\
chb04 & 0.788 & chb16 & 0.900 \\
chb05 & 0.920 & chb17 & 0.183 \\
chb06 & 0.585 & chb18 & 0.950 \\
chb07 & 0.917 & chb19 & 0.500 \\
chb08 & 0.480 & chb20 & 0.988 \\
chb09 & 0.787 & chb21 & 0.700 \\
chb10 & 0.786 & chb22 & 1.000 \\
chb11 & 0.667 & chb23 & 0.750 \\
\bottomrule
\end{tabular}
\end{table}

\end{appendices}

\end{document}
